\paragraph{Finding A09 (an explicit rank-one theta point of order thirteen).}
Let
\[
 C/\mathbb F_2:\quad Y^2+X^2Y=X^5+X,\qquad P=\infty.
\]
This curve has genus two and $2$-rank one. Put
\[
 k=\mathbb F_2[u,v]/(u^2+u+1,\ v^2+u^2v+1)\simeq\mathbb F_{16},
 \qquad R=(u,v)\in C(k).
\]
Then $L=\mathcal O_C(R-P)$ has exact order $13$ and lies on the
Raynaud theta divisor. Since $2$ has order $12$ modulo $13$,
this is a positive-$2$-rank instance of $T(C)$.

\begin{proof}
The substitution $x=X^{-1}$, $y=Y/X^2$ gives
$y^2+y=x^3+x^{-1}$. Its two Artin--Schreier poles have odd orders
$3$ and $1$. Riemann--Hurwitz gives $2g-2=-4+4+2=2$, and
Deuring--Shafarevich gives $f_C-1=2(-1)+2=0$, so $g=2$ and $f_C=1$.
There is one point $P$ above $X=\infty$, with
$\operatorname{ord}_P(X)=-2$ and $\operatorname{ord}_P(Y)=-5$.
The differential $dX/X^2=d(1/X)$ has divisor $2P$.
It is locally exact, and hence its corresponding section of the
rank-one sheaf $B^1$ has divisor $P$. Thus
$B^1\simeq\mathcal O_C(P)$ and the theta support is $C-P$.
Here the Frobenius twist is identified with the same equation over
$\mathbb F_2$; no choice of a twist over a nonperfect field is used.

The field presentation is explicit. Over $\mathbb F_4=\mathbb F_2(u)$,
putting $w=v/u^2$ gives $w^2+w=u^2$. Its trace to $\mathbb F_2$ is
$u^2+u=1$, so the second quadratic is irreducible. The relation
$u^3=1$ gives $u^5+u=1$, proving $R\in C(k)$.

Define
\begin{align*}
 A(X)={}&u(v+1)(X^6+X^3)+u^2vX^5+uX^4
                 +(uv+1)X^2+uX+u^2,\\
 B(X)={}&X^4+uX^3+X+1,\qquad F=A(X)+B(X)Y.
\end{align*}
Direct polynomial reduction in the displayed field gives
\[
 A(X)^2+X^2A(X)B(X)+(X^5+X)B(X)^2=(X+u)^{13},
 \tag{PR.1}
\]
and
\[
 F(R)=0,\qquad F(\iota R)=u\ne0,
 \qquad\iota(X,Y)=(X,Y+X^2).
 \tag{PR.2}
\]
The left side of (PR.1) is the norm $F\,\iota(F)$.
As $u\ne0$, the two points over $X=u$ are unramified and distinct.
Equations (PR.1)--(PR.2) imply that every finite zero of $F$ is $R$
and that its multiplicity is exactly $13$. There are no finite poles.
At $P$, the term $X^4Y$ has pole order $13$, while every other term
has smaller pole order. Therefore
\[
 \operatorname{div}(F)=13R-13P.\tag{PR.3}
\]
The class $R-P$ is nonzero: otherwise a rational function with that
divisor would give a degree-one map from a genus-two curve to
$\mathbb P^1$. Since $13$ is prime, (PR.3) proves its exact order.
Also $B^1\otimes L\simeq\mathcal O_C(R)$, so $L\in\Theta_{B^1}$.
Finally, the successive powers of $2$ modulo $13$ are
\[
 1,2,4,8,3,6,12,11,9,5,10,7,
\]
which are all the nonzero residues.
\end{proof}

\paragraph{An explicit cover and the arithmetic check.}
The normalization of the Kummer equation
\[
 W^{13}=F
\]
is a geometrically connected etale $\mu_{13}$-torsor over $C_k$.
Indeed, (PR.3) makes every valuation divisible by $13$, so it is
etale; it is geometrically connected because $F$ cannot be a
thirteenth power, again by nontriviality of $R-P$. With the local
parameter $t=X^2/Y$ at $P$, the unit $t^{13}F$ has residue $1$.
The root $tW=1$ therefore gives a $k$-rational point above $P$.
After extending constants to $\mathbb F_{4096}$, which contains
$k$ and $\mu_{13}$, its deck group is the constant group
$\mathbb Z/13$. Its genus is $14$ and its new Jacobian factor has
$2$-rank zero by A11 below (or by the prime case A03).
For this assertion about the displayed Kummer model, note that
$C$ and its theta support are defined over $\mathbb F_2$, so
Frobenius relabelling preserves the theta incidence; the explicit
twist normalization in A11 applies to its defining torsion class.
The direct torsor construction here is not asserted to be the same
arithmetic model as the Lang cover over $\mathbb F_{64}$ in the
recovered A1 row B.39, working identifier B.LIT7.

The file \texttt{theta13\_direct.py} constructs the identity using
Hasse--Taylor linear algebra over
$\mathbb F_2[z]/(z^4+z+1)$, with $u=z^2+z$, $v=z^3$.
Its certificate lists every coefficient. The separate verifier
\texttt{verify\_theta13\_norm.py} imports no producer code and
checks (PR.1)--(PR.2) by polynomial reduction over $\mathbb F_2$.
These are exact finite computations; they do not extrapolate to a
classification of positive-rank curves.

\paragraph{Nonabelian content of the example.}
The same curve has
\[
 \#C(\mathbb F_2)=4,\qquad \#C(\mathbb F_4)=4,\qquad
 H_2(U)=U^4+U^3+2U+4,\qquad \#J_C(\mathbb F_2)=8.
 \tag{PR.10}
\]
Indeed, the Artin--Schreier model has its two rational branch
points, and $x=1$ contributes two points. For
$x\in\mathbb F_4\setminus\mathbb F_2$, its right side is
$x^3+x^{-1}=1+x^2=x$,
whose absolute trace is $1$, so no additional points occur.
The two counts and the genus-two functional equation give the
displayed characteristic polynomial; evaluating it at $1$ gives
$8$. Thus all four rational points have the same prime-to-$2$
abelian Kummer class, because their divisor differences lie in
the $2$-group $J_C(\mathbb F_2)$. Nevertheless A09 and A11 show
that their full geometrically prime-to-$2$ point-section classes
are pairwise distinct. This is a concrete positive-rank instance
where the nonabelian covering argument supplies information
absent from the abelian quotient.

\paragraph{Finding A10 (three is excluded for every nonordinary genus-two base).}
Let $C$ be a smooth proper genus-two curve over an algebraically
closed field of characteristic two, with $f_C<2$. Then its Raynaud
theta divisor contains no nonzero point of order $3$.
Consequently every connected cyclic etale degree-three cover has
ordinary new factor of dimension two. In particular, this conclusion
also holds when $f_C=1$.

\begin{proof}
Work on the Frobenius twist, temporarily denoting that curve by
$C$ as well, and write $B=B^1$ there. In
characteristic two this is a line bundle with $B^{\otimes2}\simeq
\omega_C$ and degree one. The square relation follows in an etale
coordinate by comparing transition functions of locally exact
differentials with those of $\omega_C$.
From the relative Frobenius exact sequence,
$H^0(B)$ is the kernel of Frobenius on $H^1(\mathcal O_C)$.
Nonordinarity implies that Frobenius is not bijective, so $H^0(B)\ne0$.
Thus $B\simeq\mathcal O(W)$ for a point $W$, and
$2W\sim K_C$. Its theta support is $C-W$.

If $L\ne0$ had order $3$ on this support, write
$L=\mathcal O(R-W)$ with $R\ne W$. Then $3(R-W)$ would be
principal, so there would be a function with unique pole of exact
order $3$ at $W$. But Riemann--Roch and $2W\sim K_C$ give
\[
 h^0(2W)=2=h^0(3W).
\]
Every function with pole divisor at most $3W$ therefore already has
pole divisor at most $2W$, a contradiction.
For a cyclic degree-three cover, every nontrivial eigenspace
$H^1(C,L^i)$ is one-dimensional. The Frobenius maps between the
two eigenspaces have zero kernel precisely when the associated
theta incidences are absent, by the same exact sequence used in
A11. They are both bijective, so the new factor is ordinary.
\end{proof}

\subsection{A prime-power bridge and the remaining order obstruction}
\label{sec:prime-power-bridge}

\paragraph{Finding A11 (a sufficient prime-power incidence condition).}
Let $C$ be a smooth proper genus-two curve in characteristic two.
Let $n=\ell^r$ with $\ell$ an odd prime and $r\ge1$. Suppose that
\[
 \langle 2,-1\rangle=(\mathbb Z/n\mathbb Z)^\times
 \tag{PR.4}
\]
and that the Raynaud theta divisor contains a class of exact order
$n$. Then over an algebraic closure there is a connected cyclic etale
degree-$n$ cover $D_n\to C$ whose primitive conductor-$n$ factor has $2$-rank zero.
Writing $D_{n/\ell}$ for its intermediate degree-$n/\ell$ cover of
$C$, this factor is precisely
\[
 \operatorname{Prym}(D_n/D_{n/\ell})
\]
up to isogeny. In the finite-field covering formulation of A1, this
single cover separates every proposed pair of distinct colliding
points. Thus prime-power incidence with (PR.4) is sufficient; prime
order and the stronger equality $\langle2\rangle=(\mathbb Z/n)^\times$
are not required.

\begin{proof}
\emph{1. Symmetry and a zero Frobenius map.}
Let $B$ be the locally exact differential bundle on $C^{(1)}$ and
$L\in\operatorname{Pic}^0(C^{(1)})$ the given class. Since
$B^{\otimes2}\simeq\omega_{C^{(1)}}$ and $\chi(B\otimes L)=0$,
Serre duality gives
\[
 h^0(B\otimes L)=h^0(B\otimes L^{-1}).
 \tag{PR.5}
\]
Thus both $L$ and $L^{-1}$ lie on theta. Tensoring the relative
Frobenius sequence by $L$ gives
\[
 0\longrightarrow L\longrightarrow F_*F^*L
       \longrightarrow B\otimes L\longrightarrow0.
 \tag{PR.6}
\]
The two degree-zero line bundles $L$ and $F^*L$ are nontrivial:
$n$ is odd and Frobenius pullback preserves its exact order.
They have no sections, and each has one-dimensional $H^1$ by
Riemann--Roch. Therefore $H^0(B\otimes L)\ne0$ says that the
induced map from $H^1(L)$ to $H^1(F^*L)$ has nonzero kernel;
being a map of lines, it is zero. The same applies to $L^{-1}$.
To make the character labels precise, choose $e$ with $2e\equiv1
\pmod n$ and put $M=(F^*L)^e$ on $C$. This class still has exact
order $n$, and $M^2=F^*L$. If $M^{(1)}$ denotes its Frobenius
twist on $C^{(1)}$, then $F^*(M^{(1)})=M^2=F^*L$.
The pullback $F^*$ is injective on odd-order torsion, so
$M^{(1)}=L$. Consequently (PR.6) gives literally the
Frobenius-semilinear arrow from $H^1(C,M)$ to $H^1(C,M^2)$;
the inverse class gives the arrow from $H^1(C,M^{-1})$ to
$H^1(C,M^{-2})$.

\emph{2. Every primitive orbit is nilpotent.}
Over an algebraic closure, the Kummer correspondence gives the
connected cyclic etale cover associated to $M$, choosing its
character labels so that the eigensheaves are $M^i$.
When $C$ is over a finite field, all these data descend after a
finite constant extension. Its cohomology decomposes as
\[
 H^1(D_n,\mathcal O_{D_n})=
   H^1(C,\mathcal O_C)\oplus
   \bigoplus_{i=1}^{n-1} V_i,\qquad
   V_i=H^1(C,M^i),\quad \dim V_i=1.
 \tag{PR.7}
\]
Absolute Frobenius sends $V_i$ to $V_{2i}$. The primitive part is
the sum over $i\in(\mathbb Z/n)^\times$.
Condition (PR.4) means that the action of multiplication by $2$
on this unit group has at most two orbits, represented by $1$ and
$-1$. By Step~1, each orbit contains a zero arrow. A sufficiently
large iterate of Frobenius is therefore zero on every $V_i$ in
the primitive sum: traverse the finite orbit until the zero arrow
is reached. The stable Frobenius part of that sum is zero, which
is the assertion that the primitive factor has $2$-rank zero.
The one-dimensional spaces are those on the original genus-two
base; no inference from a single singular higher-dimensional
matrix is being made.

\emph{3. Prime powers identify the entire relative factor.}
Choose a generator $\sigma$ of the order-$n$ deck group and put
$a=\sigma^{n/\ell}$, of order $\ell$.
The quotient by $\langle a\rangle$ is $D_{n/\ell}$.
On the character with index $i$, the element $a$ acts trivially
if and only if $\ell\mid i$. For $n=\ell^r$, its nontrivial
characters are exactly the primitive indices $\gcd(i,n)=1$.
Thus the primitive factor is exactly the new factor for this
degree-$\ell$ quotient, proving the stated Prym identification.
The dimension is $(n-n/\ell)(g(C)-1)=\varphi(n)$, as a check.

\emph{4. The point-separation consequence.}
Now work in the finite-field situation of the covering formulation
of A1. Enlarge constants to define the cover, make its deck group
constant, and choose a rational point $P_n$ above the given base
point $P$. If a distinct point $Q$ collided with $P$, the covering
condition would provide a rational $Q_n$ above $Q$ with
\[
 \delta=[Q_n-P_n]\in J_{D_n}[2^\infty].
\]
The point $(a-1)\delta$ is still $2$-primary and belongs to
$\operatorname{im}(a-1)=\ker(\operatorname{Nm})^0
=\operatorname{Prym}(D_n/D_{n/\ell})$.
An abelian variety of $2$-rank zero has no nonzero geometric
$2$-power torsion points. Thus $(a-1)\delta=0$.
Equivalently, the divisor
\[
 aQ_n+P_n-Q_n-aP_n
 \tag{PR.8}
\]
is principal. Its positive and negative parts both have degree two
and are disjoint: the deck action is free, and the orbits above
the distinct base points $P$ and $Q$ are disjoint. It therefore
defines a degree-two map $D_n\to\mathbb P^1$.

Such a map is separable, since a purely inseparable map of degree
two from a smooth curve over an algebraically closed field to
$\mathbb P^1$ would make the curve rational. Hence $D_n$ would be
hyperelliptic. Its genus is $n+1\ge4$, so its hyperelliptic pencil
is unique. The order-$\ell$ automorphism $a$ acts on its base
$\mathbb P^1$; this action is nontrivial, because the kernel of
the action on a separable degree-two function field has order
at most two. A nontrivial odd-order automorphism of $\mathbb P^1$
has a fixed point. The fibre over that point contains at most two
geometric points, and an odd-prime cyclic group acts trivially on
a set of size at most two. This supplies a fixed point of $a$ on
$D_n$, contradicting the freeness of the etale deck action.
Thus the proposed collision is impossible.
\end{proof}

Condition (PR.4) is a finite group computation for the actual integer
$n$ in question. No non-Wieferich assumption is needed for the
statement or proof. Hooley's theorem supplies prime-level
primitive-root information under its analytic hypotheses; it does
not supply a theta point of any of these prime or prime-power
orders, nor does it imply (PR.4) at an arbitrary higher power.

\paragraph{Attempted extension using primary projections.}
For every genus-two curve over $\overline{\mathbb F}_2$, the
Raynaud theta support is a translate of its Abel curve, since
$\deg B=1$. It therefore generates its Jacobian: a proper translate
of a subgroup containing theta would contain all differences of
points of the Abel curve, whereas those differences generate the
Jacobian.
Poonen's Theorem~1.8 consequently applies to this theta curve:
given a finite set $S_0$ of primes, there is a density-one set
$S\supseteq S_0$ for which projection of theta points onto the
$S$-primary torsion of $J$ is onto \cite[Theorem~1.8]{Poonen2005}.
In particular, take $S_0$ to contain $2$ and a selected odd prime
$\ell$, prescribe an exact-order-$\ell^r$ point in the $\ell$-part,
and prescribe zero at every other prime of $S$. A resulting theta
point has order
\[
 n=\ell^r m,\qquad (\ell,m)=1,\qquad
 \text{every prime dividing }m\text{ lies outside }S.
 \tag{PR.9}
\]
This gives odd-order points, and iterating with all previously
occurring primes in $S_0$ gives infinitely many of pairwise coprime
orders. It applies also to ordinary curves with split Jacobians.
It does not require or prove that $m=1$.
Indeed, for a nonordinary genus-two curve take $\ell=3$ and $r=1$.
Finding A10 proves that the resulting point cannot have order $3$,
so its extra factor $m$ is necessarily greater than one. The
uncontrolled factors cannot simply be omitted from this projection
argument.

Two attempted ways to remove $m$ fail at definite algebraic steps.
First, multiplying the point by $m$ removes the extra primary
components but places it on $[m](\Theta)$; this operation does not
preserve $\Theta$. Second, pass from a cyclic degree-$\ell m$
cover to its intermediate degree-$m$ base. Its relative
degree-$\ell$ factor contains the characters of every conductor
$\ell d$ with $d\mid m$, including conductor $\ell$ itself.
Theta incidence at exact order $\ell m$ controls only the
primitive conductor-$\ell m$ block. On the intermediate base the
character spaces have dimension $m$, so one zero constituent does
not make those whole spaces nilpotent.

Even if $\langle2,-1\rangle=(\mathbb Z/\ell m)^\times$, symmetry
and that group equality kill only the primitive block. They leave
the proper-conductor blocks untouched. More intrinsically, when
$n$ has at least two distinct prime factors, the least common
multiple of its proper divisors is $n$. A deck element that acts
trivially on all proper-conductor characters must therefore be
the identity; there is no nontrivial element to substitute for
$a$ in Step~4 using only the primitive-zero conclusion. For a prime
power, that least common multiple is $n/\ell$, precisely the
reason Step~3 works.

The remaining finite-field incidence question is consequently
sharpened but unresolved here: for a positive-$2$-rank genus-two
curve, must theta contain a point of some odd prime-power order
$n$ satisfying (PR.4), or can another argument control all the
needed proper-conductor blocks? Neither the density-one projection
theorem nor the prime distribution theorem proves this assertion.
